\documentclass[11pt]{article}

\usepackage[letterpaper,margin=0.85in]{geometry}
\usepackage{mathpazo}
\usepackage[T1]{fontenc}
\usepackage[utf8]{inputenc}
\usepackage{microtype}
\usepackage{enumitem}
\usepackage{booktabs}
\usepackage{tabularx}
\usepackage{array}
\usepackage{fancyhdr}
\usepackage{hyperref}
\usepackage{listings}
\usepackage{xurl}

\hypersetup{hidelinks}
\setlength{\parindent}{1.25em}
\setlength{\parskip}{0.2em}
\setlist{nosep,leftmargin=1.4em}
\renewcommand{\arraystretch}{1.2}
\newcolumntype{Y}{>{\raggedright\arraybackslash}X}
\pagestyle{fancy}
\fancyhf{}
\lhead{API 209 Math Camp 2026}
\rhead{Lab 2: Build the analysis dataset}
\cfoot{\thepage}
\setlength{\headheight}{14pt}

\lstset{
  basicstyle=\ttfamily\footnotesize,
  breaklines=true,
  breakatwhitespace=true,
  columns=fullflexible,
  frame=single,
  framerule=0.4pt,
  xleftmargin=0.4em,
  xrightmargin=0.4em,
  aboveskip=0.5em,
  belowskip=0.5em,
  showstringspaces=false
}

\newcommand{\checkpoint}[2]{\medskip\noindent\textbf{Checkpoint #1. #2}\par}
\newcommand{\file}[1]{\nolinkurl{#1}}
\newcommand{\answerline}{\par\vspace{0.45em}\noindent\rule{\linewidth}{0.35pt}\par}

\begin{document}

\begin{center}
  {\small THURSDAY, AUGUST 20 \quad 4:00--5:00 P.M.}\par
  \vspace{0.6em}
  {\LARGE\textbf{Lab 2: From a draft file to an analysis dataset}}\par
  \vspace{0.45em}
  {\large Commission the work, interrogate the evidence, and decide what survives}
\end{center}

\section*{The policy brief}

You are preparing evidence for this question:

\begin{quote}
\textbf{How is national income associated with under-five mortality across
economies between 2000 and 2022?}
\end{quote}

The finished dataset must contain one observation per economy--year, GDP per
capita (PPP), under-five mortality, a natural-log measure of income, and the
economy information needed for interpretation. Observations without usable
income or mortality measures do not enter this first analysis.

Work from the RStudio project, this handout, and
\file{code/lab-2-starter.R}. The script is an audit instrument: it exposes
evidence and evaluates a candidate output, but it does not build the final
pipeline for you.

\subsection*{Two uses of Codex}

\noindent\begin{tabularx}{\textwidth}{@{}p{0.22\textwidth}Y Y@{}}
\toprule
 & \textbf{Codex as verifier} & \textbf{Codex as pipeline operator} \\
\midrule
Assignment & Examine the available files and make evidence-backed claims. &
Choose, implement, and run a complete construction strategy. \\
Risk & A correct calculation can support an unjustified explanation. & A
pipeline can run, save, and pass structural checks while embodying a bad
decision. \\
Your responsibility & Reproduce consequential claims and challenge unsupported
inferences. & Inspect the source choice, code, output, and tests before accepting
the result. \\
\bottomrule
\end{tabularx}

Codex output is a candidate contribution to the analysis. Fluency, successful
execution, and a passing test are evidence of different things; none settles
the substantive judgment by itself.

\newpage
\section*{Round I: Codex as verifier}

Before inspecting either file, specify the standard by which you will judge a
candidate dataset. At minimum, state the unit of observation, identifying
variables, period, required measures, exclusions, and one downstream result
that a record-level error could change.

\checkpoint{1}{State the analytical risk before seeing the records}

Describe one plausible data problem and trace its consequences through a
briefing product:

\begin{quote}
record-level problem $\longrightarrow$ changed table, plot, or estimate
$\longrightarrow$ changed policy interpretation.
\end{quote}

\answerline\answerline

Run the \textbf{evidence audit} in \file{code/lab-2-starter.R}. Use its output
to decide which observations deserve investigation and which comparisons are
needed to resolve them.

Then give Codex this verification assignment:

\begin{lstlisting}
Act as a skeptical data reviewer for a policy briefing on national income and
under-five mortality across economies from 2000 to 2022.

Read the handout, data/lab-2/health-briefing-draft.csv,
data/derived/math-camp-wdi-2000-2022.csv, and
data/documentation/indicator-dictionary.csv. Work read-only.

Assess whether the draft is fit for the stated analysis. Examine its unit of
observation, economy-year key, coverage, required variables, missingness,
ranges, and disagreements with the documented project data. For every
conclusion, report the file, variable, economy-year when relevant, and observed
value. Distinguish what the files establish from your explanation of how a
problem may have arisen. End with: (1) your recommendation, (2) the two claims
that matter most for the briefing, and (3) reproducible R checks for them.

Do not edit, deduplicate, or rebuild any file.
\end{lstlisting}

\checkpoint{2}{Adjudicate the verification report}

Select the most consequential claim in the response. Reproduce its evidence
with R, then record a verdict: \emph{accept}, \emph{revise}, or \emph{reject}.

\noindent\begin{tabularx}{\textwidth}{@{}p{0.25\textwidth}Y@{}}
\toprule
Claim & \rule{0pt}{2.0em} \\
Observed evidence & \rule{0pt}{2.0em} \\
What the agent inferred & \rule{0pt}{2.0em} \\
Verdict and consequence for the briefing & \rule{0pt}{2.0em} \\
\bottomrule
\end{tabularx}

Identify one sentence in the response that is more certain than its evidence.
If you find none, write the strongest counterargument to the recommendation.

\answerline

\newpage
\section*{Round II: Codex as pipeline operator}

Now commission the complete construction. The assignment deliberately states
the analytical requirements without prescribing every implementation choice.
The agent must make those choices visible so that you can judge them.

\begin{lstlisting}
Build and run the analysis-data pipeline for this policy question:
How is national income associated with under-five mortality across economies
between 2000 and 2022?

Read the project documentation and relevant data files. Decide which input is
authoritative and justify that decision from project evidence.

Create code/02-health-agent-pipeline.R. It must save:
  outputs/02-health-analysis.csv
  outputs/02-health-build-record.csv

Requirements for the analysis file:
- one row per iso3c + year;
- years 2000 through 2022;
- iso3c, country, region, current income classification, year, GDP per capita
  (PPP), under-five mortality, and log income;
- exclude missing or nonpositive income and missing mortality;
- calculate log income as the natural logarithm of GDP per capita (PPP).

You may create or edit only code/02-health-agent-pipeline.R and the two named
outputs. Do not modify source data, setup files, the starter audit, or the
independent verifier. Run the pipeline. Report the input chosen, exclusions,
files changed, checks performed, and any analytical decision that remains
contestable.
\end{lstlisting}

\checkpoint{3}{Review the commission before accepting execution}

What has been delegated? What decision remains yours? Name one way the agent
could satisfy the file specification yet produce a substantively indefensible
dataset.

\answerline\answerline

Inspect the new script and the agent's report. Do not evaluate the final output
only from the prose summary. Look for the chosen input, any rule that resolves
multiple records, sample restrictions, transformation, and actual files
written.

\newpage
\section*{Audit the pipeline, not merely its exit status}

Run the \textbf{candidate-output audit} in
\file{code/lab-2-starter.R}. It reports two distinct forms of evidence:

\begin{enumerate}
  \item \textbf{structural evidence}: required columns, unique keys, period,
  missingness, ranges, and recomputation of logged income; and
  \item \textbf{source concordance}: missing or unexpected keys and values that
  disagree with the documented project data.
\end{enumerate}

\checkpoint{4}{Find the strongest reason to accept or reject the pipeline}

\noindent\begin{tabularx}{\textwidth}{@{}p{0.31\textwidth}Y@{}}
\toprule
Does the script implement the stated unit and sample? & \rule{0pt}{2.2em} \\
Does the output pass structural checks? & \rule{0pt}{2.2em} \\
Does the output agree with documented source records? & \rule{0pt}{2.2em} \\
Is every conflict-resolution rule defensible without relying on row order? & \rule{0pt}{2.2em} \\
Decision: accept, revise, or reject & \rule{0pt}{2.2em} \\
\bottomrule
\end{tabularx}

A pipeline can pass one set of checks and fail another. If everything passes,
identify a consequential question that the current audit still cannot answer.

\newpage
\section*{From checked data to a briefing table}

Use the candidate only after you have accepted it or revised the pipeline. A
policy colleague now asks:

\begin{quote}
\textbf{In 2022, how did typical income and under-five mortality differ across
current World Bank income groups?}
\end{quote}

Write an R pipeline that starts from
\file{outputs/02-health-analysis.csv}, restricts the data to 2022, groups by
\texttt{income\_level\_current}, and creates one row per income group with:

\begin{itemize}
  \item the number of economies;
  \item median GDP per capita (PPP); and
  \item median under-five mortality.
\end{itemize}

Arrange the table from lowest to highest median mortality and save it as
\file{outputs/02-health-briefing-by-income.csv}. This is an unweighted
descriptive comparison: each economy contributes once to its group's 2022
summary.

\checkpoint{5}{Defend the grouped briefing result}

Before computing, state how the unit changes from the analysis file to the
briefing table and how a repeated economy--year could alter a group median or
count. After computing, confirm that the group counts sum to the number of
distinct economies observed in 2022.

Write two briefing sentences: one numerical comparison supported by the table
and one limitation that prevents a causal or historical-classification claim.

\answerline\answerline\answerline

Ask Codex to verify the grouped result without rewriting it: independently
recompute the counts and medians from the accepted analysis file, compare them
with the saved briefing table, and identify any sentence that goes beyond the
evidence.

\checkpoint{6}{Distinguish the two agent roles}

In no more than four sentences, explain:

\begin{enumerate}
  \item what Codex contributed as a verifier;
  \item what it contributed as a pipeline operator;
  \item one mistake or overclaim you caught---or the most credible failure you
  attempted to falsify; and
  \item whether you would release the dataset and briefing table, and why.
\end{enumerate}

\answerline\answerline\answerline

Keep the prompt, generated pipeline, audit evidence, and your decision. The
reusable skill is not following a prescribed sequence. It is commissioning
analytical work precisely and knowing what evidence would make you reject it.

\end{document}
